\begin{code}[hide]
module slides.popl.sections.catalogue where

open import Algebra.Bundles using (Group ; Monoid)
open import Algebra.Morphism.Structures using (module MonoidMorphisms)
open import Data.Nat using (ℕ ; suc ; 2+)
open import Data.Nat.Primality using (Prime)
open import Level using (0ℓ)

open import Word.Base using (_ʷ)
open import Presentation.Definitions using (_IsPresentationOf_)
open import Presentation.Construct.Base using (_⋄_⋄_ ; CommRel ; ConjRelʷ ; _⊕^_)
import Presentation.Properties as PP
open import Examples.Groups.Cyclic.Normalization using (_Cn,_===_)
open import Examples.Groups.Symmetric.Syntactics using (_VRel,_===_)
import Examples.Groups.Pauli.Presentation as Pauli
import Examples.Construct.SemiDirectProduct.SnD as SnD
import Examples.Amalgamations.CliffordT1 as CliffordT1

open SnD using (conj ; module Wreath)
open Pauli using (Pauli-group) renaming (Γ-H to Γ-Pauli₁)
open Monoid using (rawMonoid)
open MonoidMorphisms using (IsMonoidIsomorphism)
open PP using (•-ε-monoid)
open CliffordT1.CliffordT1 using (_===_ ; amalPres ; f)
\end{code}
%% The catalogue (paper §5).
\section{A catalogue of verified presentations}

\begin{frame}{What the library has verified}
\centering
\footnotesize
\setlength{\tabcolsep}{4pt}
\begin{tabular}{@{}lll@{}}
\toprule
relation set & presents & what it exercises\\
\midrule
$T^{\,n+1}=\varepsilon$\;/\;$T^{\,0}=\varepsilon$ & $\mathbb{Z}_{n+1}$\;/\;the \emph{monoid} $(\mathbb{N},+,0)$ & monoid vs.\ group presentations\\
\aid{order}, \aid{yang-baxter} + structure & permutations of $\mathsf{Fin}\;n$ & circuits, a recursive tower\\
$(\mathbb{Z}_N)^n \rtimes S_n$-circuits & $\mathbb{Z}_N\wr S_n$, every $N,n$ & products with semantics\\
$(\mathbb{Z}_p\times\mathbb{Z}_p)^n$ & Pauli group mod phases, every prime $p$ & pure composition\\
qubit Clifford+$T$ rules & $\cong\langle T,M\rangle\ast_M\langle H,M\rangle$ & amalgamation, \aid{refl}\\
qutrit Clifford+$T$ rules & $\cong$ an amalgam over $M=\langle S,X,\zeta,H^2\rangle$ & long towers, twisted relations\\
$U_3(\mathbb{Z}[\frac12,i])$, S1--S14 & $\cong(\mathbb{Z}_4\wr S_3)\ast_M\langle M,K_{01}\rangle$ & reuse of the wreath product\\
\bottomrule
\end{tabular}

\vspace{10pt}
\begin{columns}[c]
\begin{column}{0.45\textwidth}
\begin{win}
\small Every theorem restated with its full type in one index module,
\texttt{MainTheorems.agda}.
\end{win}
\end{column}
\begin{column}{0.45\textwidth}
\small $T^0=\varepsilon$ presents $\mathbb{Z}$ as a \emph{group} but $\mathbb{N}$ as a
\emph{monoid}: words have no formal inverses, and \aid{Grouplike} is exactly what fails.
\end{column}
\end{columns}
\end{frame}

\begin{frame}{Whole families at once}
\renewcommand{\AgdaCodeStyle}{\footnotesize}
\begin{columns}[c]
\begin{column}{0.58\textwidth}
\textbf{Wreath products}: signed permutations at $N=2$, monomial matrices with
$N$-th roots of unity in general --- completeness over \alert{two naturals}:
\begin{code}
wreath-presentation : ∀ n m →
  (((suc m Cn,_===_) ⊕^ n) ⋄ (n VRel,_===_) ⋄ ConjRelʷ conj)
    IsPresentationOf (Wreath.wreath-group n m)
\end{code}
\begin{code}[hide]
wreath-presentation n m = Wreath.presentation n m
\end{code}
\vspace{4pt}
\textbf{Pauli groups} for every prime: three construction theorems composed,
\alert{no coset enumeration at all}:
\begin{code}
pauli-presentation : ∀ m (prime : Prime (2+ m)) n →
  (Γ-Pauli₁ m prime ⊕^ n) IsPresentationOf (Pauli-group m prime n)
\end{code}
\begin{code}[hide]
pauli-presentation m prime = Pauli.Pauli-presentation m prime
\end{code}
\end{column}
\begin{column}{0.39\textwidth}
\centering
{\small normal form of $\mathbb{Z}_N\wr S_n$:}\\[4pt]
\begin{tikzpicture}[x=4.6mm,y=5.5mm,line width=0.7pt]
  \foreach \y in {0,1,2,3} \draw[wire] (0,\y) -- (11,\y);
  \foreach \y/\e in {0/2,1/0,2/3,3/1}
    \node[draw,fill=tint,font=\scriptsize,minimum width=6mm,minimum height=4mm,inner sep=1pt] at (1.6,\y) {$T^{\e}$};
  \foreach \x/\p in {4/2,5.5/1,7/2,8.5/0,10/1} {
    \fill[white] (\x-0.55,\p-0.05) rectangle (\x+0.55,\p+1.05);
    \draw[wire] (\x-0.55,\p) -- (\x+0.55,\p+1);
    \draw[wire] (\x-0.55,\p+1) -- (\x+0.55,\p);
  }
  \draw[decorate,decoration={brace,amplitude=3pt,mirror}] (0.6,-0.6) -- (2.6,-0.6) node[midway,below=2pt,font=\scriptsize]{residues};
  \draw[decorate,decoration={brace,amplitude=3pt,mirror}] (3.3,-0.6) -- (10.7,-0.6) node[midway,below=2pt,font=\scriptsize]{staircases};
\end{tikzpicture}

\vspace{10pt}
{\small normal form of the Pauli group: $X^{a_1}Z^{b_1}\cdots X^{a_n}Z^{b_n}$ ---
nothing proved about Paulis specifically}
\end{column}
\end{columns}
\end{frame}

\begin{frame}{Amalgamation rediscovers Matsumoto--Amano}
\begin{columns}[c]
\begin{column}{0.52\textwidth}
\footnotesize Single-qubit Clifford+$T$; generators $T,H,S,\omega$, $X := HSSH$;
$\omega^8 = S^4 = H^2 = \varepsilon$, $T^2 = S$, $(SH)^3 = (TX)^2 = \omega$,
$\omega$ central.
\small
\[
  \text{Clifford}{+}T \;\cong\; \langle T, M\rangle \ast_M \langle H, M\rangle,
  \qquad M = \langle X, S, \omega\rangle
\]
The alternating normal form, read right to left:
\[
  (\varepsilon\mid T)\;(HT\mid SHT)^{*}\;(\varepsilon\mid H\mid SH)\;\mathcal{M}
\]
--- \alert{Matsumoto--Amano}, from the group theory alone.
\renewcommand{\AgdaCodeStyle}{\scriptsize}
\begin{code}
clifford+T-qubit-isomorphism : IsMonoidIsomorphism
  (rawMonoid (•-ε-monoid _===_)) (rawMonoid (•-ε-monoid amalPres)) (f ʷ)
\end{code}
\begin{code}[hide]
clifford+T-qubit-isomorphism = CliffordT1.CliffordT1.CliffordT1-isomorphism
\end{code}
\end{column}
\begin{column}{0.45\textwidth}
\centering
\begin{tikzpicture}[>={Stealth[length=1.8mm]},
  lab/.style={font=\tiny,align=center,text=hot}]
  \node[bx=prim,font=\scriptsize] (b) at (0,0) {$\langle\omega\rangle\times\langle S\rangle$};
  \node[bx=prim,font=\scriptsize] (m) at (0,1.1) {$M=\langle X,S,\omega\rangle$};
  \node[bx=prim,font=\scriptsize] (t) at (-1.6,2.2) {$\langle T,M\rangle$};
  \node[bx=prim,font=\scriptsize] (h) at (1.6,2.2) {$\langle H,M\rangle$};
  \node[bx=ok,font=\scriptsize] (a) at (0,3.3) {$\langle T,M\rangle\ast_M\langle H,M\rangle$};
  \draw[->,thick] (b) -- node[right,lab]{2 cosets} (m);
  \draw[->,thick] (m) -- node[left,lab]{2 cosets} (t);
  \draw[->,thick] (m) -- node[right,lab]{3 cosets} (h);
  \draw[->,thick] (t) -- (a); \draw[->,thick] (h) -- (a);
\end{tikzpicture}

\vspace{3pt}
\begin{tikzpicture}[x=7.2mm,
  s/.style={draw,minimum height=5.5mm,minimum width=6.8mm,font=\scriptsize,inner sep=1pt}]
  \node[s,fill=gray!15] at (0,0) {$\mathcal{M}$};
  \node[s,fill=tintok] at (1,0) {$HS$};
  \node[s,fill=tint] at (2,0) {$T$}; \node[s,fill=tintok] at (3,0) {$H$};
  \node[s,fill=tint] at (4,0) {$T$}; \node[s,fill=tintok] at (5,0) {$HS$};
  \node[s,fill=tint] at (6,0) {$T$};
\end{tikzpicture}

{\scriptsize a normal form with $T$-count 3}
\end{column}
\end{columns}
\vspace{4pt}
{\small Same story for qutrit Clifford+$T$ (Glaudell--Ross--Taylor's form) and for
$U_3(\mathbb{Z}[\frac12,i])$ --- where the regular expression is a \alert{new}
canonical form.}
\end{frame}
